\scanpage{256}
% archive sheet number: 255
diagramme de suites exactes
\diagram{\begin{tikzcd}[ampersand replacement=\&,column sep=small,row sep=small,cells={font=\scriptsize}]
{} \& H^i(X_0,\mathbb Z)\arrow[r]\arrow[d,"u^i"] \& R^if_*(\mathcal O_X)_y\arrow[r]\arrow[d,"v^i"] \& R^if_*(\mathcal O_X^\times)_y\arrow[r,"\partial^i"]\arrow[d,"w^i"] \& H^{i+1}(X_0,\mathbb Z)\arrow[r,"\alpha^{i+1}"]\arrow[d,"u^{i+1}"] \& R^{i+1}f_*(\mathcal O_X)_y^{\unsure{\times|\vphantom{x}}}\arrow[d,"v^{i+1}"]\\
(D)\ \biggl\{ \& H^i(X_0,\mathbb Z)\arrow[r,dash]\arrow[d] \& R^if_*(\mathcal O_X)_y^{\wedge}\arrow[r,dash]\arrow[d] \& \underset{\leftarrow}{\lim}H^i(X_n,\mathcal O_{X_n}^{\times})\arrow[r,dash,"\widehat{\partial}^{i}"]\arrow[d] \& H^{i+1}(X_0,\mathbb Z)\arrow[r,"\widehat\alpha^{i+1}"]\arrow[d] \& R^{i+1}f_*(\mathcal O_X)_y^{\wedge}\\
\deleted{\unsure{0|o}}\arrow[r,dash] \& \deleted{H^i(X_0,\mathbb Z)}\arrow[r] \& \deleted{H^i(X_n,\mathcal O_{X_n})_{\unsure{n\in\mathbb N|n\ge0}}}\arrow[r] \& \deleted{H^i(X_n,\mathcal O_{X_n}^*)_{\unsure{n\in\mathbb N|n\ge0}}}\arrow[r] \& \deleted{H^{i+1}(X_0,\mathbb Z)}\arrow[r] \& \deleted{\unsure{R|H}^{i+1}(X_n,\mathcal O_{X_n})_{\unsure{n\in\mathbb N|n}}}
\end{tikzcd}}

On en conclut, grâce au lemme des $\overline{5}$

\uline{Corollaire} 1.6.\quad a) $R^if_*(\mathcal O_X^\times)_y\to\underset{\leftarrow n}{\lim} H^i(X_n,\mathcal O_{X_n}^\times)$ est \uline{injectif}, et les deux mbs ont $=$ image ds $H^i(X_n,\mathcal O_{X_n}^*)$.

b) \deleted{\unsure{Supposons}} Pour que $R^if_*(\mathcal O_X)_y$ soit \unsure{artinien}, i.e. $y$ un pt isolé de $\operatorname{supp}R^if_*(\mathcal O_X)$. Alors \deleted{\unsure{Alors}} \inserted{$\underset{U\ni y}{\underrightarrow{\lim}}H^i(f^{-1}(U),\uline{\mathcal O}_U^*)={}$} $R^if_*(\uline{\mathcal O}_X^\times)\to\underset{\leftarrow}{\lim}H^i(X_n,\uline{\mathcal O}_{X_n}^\times)$ \deleted{soit} un \uline{isomorphisme}.
% margin note written sideways
\margin{\deleted{\unsure{Comme les systèmes projectifs des} \illegible{1}}}
% slanted
\margin{A \unsure{expliciter|contre-exples} : pour $n$ grand Néron-Severi $(X_0)$ \unsure{reste constant} c'est i.e. :
$\operatorname{Ker}(H^2(X_0,\mathbb Z)\longrightarrow R^2f_*(\mathcal O_X)_{\unsure{y|0}})$}

c) \deleted{\unsure{Proposition} \unsure{1.7}.} Le lemme des $\overline{5}$ montre également que $R^if_*(\uline{\mathcal O}_X^\times)\longrightarrow\underset{\leftarrow}{\lim}H^i(X_n,\mathcal O_{X_n}^\times)$ est toujours un monomorphisme.
\deleted{\unsure{Comme|Si} $\unsure{\beta|\delta|\xi}=0$, $\{\underset{\leftarrow}{\lim}H^i(X_n,\mathcal O_{X_n}^\times),i\}$, il existe un \unsure{syst. proj.} \illegible{2} $(\unsure{R|H}^i(X_n,\mathcal O_{X_n}^*))_n$ satisfait M.L. (A.R. \unsure{ceci} \illegible{1})}
\deleted{\unsure{Comme} $v^{i+1}$ est injectif, \unsure{on} $[\,\operatorname{Ker}\alpha^{i+1}=\operatorname{Ker}\widehat\alpha^{i+1}\,]$, donc $\operatorname{Im}\partial^i=\operatorname{Im}\widehat{\partial}^i$ donc pour}
% slanted
\margin{$g\in R^if_*(\mathcal O_X)_y$ ayant \unsure{$\overline{m}$ image|m. image|$=$ image} dans $H^{i+1}(X_0,\mathbb Z)$ on \unsure{conclut} que \unsure{$Dg$ et $g$} sont congrus mod $\operatorname{Im}H^i(X_{\unsure{U|n|0}},\mathcal O_{X_{\unsure{y|n}}})$,
\uuline{Or} $R^if_*(\mathcal O_{X_U})\to H^i(X_0,\mathcal O_{X_n})$ est surjectif, on peut corriger, O.K.}

Le cas le plus important pour nous est celui où \unsure{$\nu$|$i$|$v$}$=1$, qui donne le

\uline{Th.} 1.\deleted{\unsure{8}}7. Soit $f:X\to Y$ un morphisme propre d'espaces analytiques \uline{sur} $\mathbb C$, $y\in Y$,
\inserted{$Y_n=\operatorname{Specan}\mathcal O_y/\mathfrak m_y^{n+1}$,} $X_n=X\times_Y Y_n$, \deleted{\unsure{Alors}} posons
\[
\begin{aligned}
\operatorname{Pic}(X_y)&=\underset{U\ni y}{\underrightarrow{\lim}}\operatorname{Pic}(f^{-1}(U))\\
&=R^1f_*(\uline{\mathcal O}_X^\times)_y,\quad \deleted{\text{Alors considérons l'hom.}}\\
\operatorname{Pic}(\widehat X_y)&=\underset{\leftarrow n}{\lim}\operatorname{Pic}(X_n).
\end{aligned}
\]
